\section{Introduction}\label{sec:introduction}

Zero-shot audio classification uses class semantics to recognize sound classes without training examples from the target task, reducing the need for annotated audio. Contrastive Language--Audio Pretraining (CLAP) learns a shared representation space through audio--text contrastive learning, allowing natural-language descriptions to define classes for sound classification~\cite{elizalde2023clap}. However, short class names provide limited semantic information. Sound sources, contexts, and relations to other events are often not fully expressed in the text input.

Prior studies have enriched class semantics through textual prompting and structured knowledge. ReCLAP rewrites training captions and generates custom class prompts to add acoustic and contextual information~\cite{ghosh2024reclap}. \redrevision{PAT uses audio from the downstream task to assign weights to prompts and combine them. It also aligns audio and text representations~\cite{seth2025pat}.} These approaches improve the use of class semantics through richer text and better representation alignment. Knowledge graphs explicitly represent associations between sound events. iKnow-audio constructs an Audio-centric Knowledge Graph (AKG) and retrieves tail entities associated with candidate classes through a curated set of relations. It concatenates the tail entities with class names to form \greenrevision{enriched prompts} and adjusts CLAP's class rankings~\cite{olvera-etal-2025-iknow}. This study shows that graph knowledge associated with classes can support audio classification. \redrevision{Using this knowledge for a given audio sample also requires assessing its discriminative value for that sample.}

Using graph knowledge for classification involves relation selection, textual expression, and score fusion. First, graph retrieval identifies knowledge associated with a candidate class, but this association alone does not establish its discriminative value for the current audio. The value of a relation may vary across audio samples, motivating relation selection based on predictions for the current input. Second, concatenating a class name and a tail entity retains both entities but leaves their connecting relation implicit. The resulting text therefore does not fully express the retrieved relational information. Text construction should preserve the relation's meaning while limiting additions unsupported by the triple. Finally, knowledge-augmented prompts provide additional evidence, but differ in how closely they match the current audio. \redrevision{Knowledge weakly related to the input audio may also alter candidate class rankings.} Knowledge augmentation therefore needs to retain the discriminative information in the original CLAP scores while using graph evidence to refine predictions.

We propose Sample-Adaptive Knowledge Integration (SAKI), which converts previously retrieved class-related knowledge into evidence for individual audio samples while keeping CLAP parameters unchanged. SAKI constructs knowledge text offline, then selects relations for the input audio during online inference. It combines knowledge scores with the original CLAP scores to rerank candidate classes.

The main contributions of this work are as follows:
\begin{itemize}
  \item We introduce sample-level relation selection using top-1 voting and top-two score gaps from the class predictions obtained for each relation. This selects relation sets for individual audio samples without additional training.
  \item We introduce Audio-Aligned Knowledge Verbalization (AAKV), which converts knowledge graph triples into descriptions of sounds while preserving relations between entities. These descriptions express relational information as text for audio--text matching.
  \item We develop Hierarchical Fusion to aggregate knowledge evidence within each class and combine the aggregated score with the original CLAP score through an explicit weight. \redrevision{This separates evidence aggregation from the weighting of knowledge evidence in the final class scores.}
\end{itemize}

We evaluate SAKI on five public audio datasets against our controlled reimplementation of iKnow-audio (iKnow\textsuperscript{\dag}). SAKI improves Hit@1 by 2.08--12.73 percentage points and mean reciprocal rank (MRR) by 1.02--6.68 percentage points. In the component ablation, SAKI achieves the highest unweighted mean Hit@1 and MRR across the five datasets. Comparisons of relation selection and knowledge verbalization further support the respective designs under the evaluated settings.

The remainder of this paper is organized as follows. Section~2 reviews related work, and Section~3 presents the SAKI framework and its core components. Section~4 reports the experimental setup, results, analyses, and limitations. Section~5 concludes the paper.
